problem:  
Let \( (G/H,g) \) be a compact connected Riemannian homogeneous space with reductive decomposition \( \mathfrak g = \mathfrak h \oplus \mathfrak m \) and \( g \) induced by an \( \mathrm{Ad}(H) \)-invariant inner product on \( \mathfrak m \).  

Define the **set of geodesic vectors in \(\mathfrak g\)** by  
\[
\mathcal{G} := \{ X\in \mathfrak g \mid \langle [X, Y]_{\mathfrak m}, X_{\mathfrak m}\rangle = 0,\ \forall Y\in \mathfrak m \},
\]
and let \(\pi_{\mathfrak m} : \mathfrak g \to \mathfrak m\) be the canonical projection, so the **geodesic‑direction set** is \(\mathcal{D}_o := \pi_{\mathfrak m}(\mathcal{G})\).  

**Open problem (precise form):**  
Assume \( (G/H,g) \) is compact and geodesically complete. Suppose  
1. the set \(\mathcal{G}\subset \mathfrak g\) is a smooth \(\mathrm{Ad}(H)\)-invariant submanifold of \(\mathfrak g\);  
2. for each \(X\in \mathcal{G}\), the one‑parameter subgroup \(\exp(tX)\) is periodic (i.e. there exists \(T_X>0\) such that \(\exp(T_X X)\in H\)).  

**Conjecture:**  
Under these conditions, must \( (G/H,g) \) be **naturally reductive**?  
Equivalently, does smoothness of the set of geodesic vectors together with periodicity of all corresponding homogeneous geodesics force the Lie algebraic identities  
\[
\langle [Z, X]_{\mathfrak m}, Y\rangle + \langle X, [Z,Y]_{\mathfrak m}\rangle = 0,\quad \forall X,Y,Z\in \mathfrak m\,?
\]
If not, can there exist non‑naturally‑reductive compact homogeneous spaces for which \(\mathcal{G}\) is a smooth submanifold and all homogeneous geodesics are closed?

Why is it a "good" problem:  
This refinement closes the earlier logical gaps by shifting the focus to \( \mathcal G\subset \mathfrak g\), where the geodesic lemma is naturally formulated, and using the one‑parameter subgroup directly for the periodicity condition. The problem now asks whether **algebraic regularity of the space of geodesic vectors** implies global geometric symmetry.  

The question is deceptively simple: if the algebraic variety defined by the geodesic lemma happens to be smooth, does that enforce the structure equations of a naturally reductive space? Behind this stand deep, intertwined categories—**Lie algebraic geometry** (smoothness of a constraint variety), **homogeneous geodesic dynamics** (periodicity of exponentials), and **metric geometry** (natural reductivity). The problem isolates a new possible *regularity‑rigidity* principle: smoothness of the geodesic‑vector manifold might act as an algebraic signature of natural reductivity.  

Known non‑naturally‑reductive g.o. examples typically yield singular, stratified sets of geodesic vectors; proving or disproving this conjecture would uncover whether this singularity is essential. Thus the problem leverages a clean and well‑posed condition to connect **algebraic geometry of the Lie algebra** with **Riemannian homogeneous rigidity**, offering a clear, well‑defined, and potentially far‑reaching path for new research.